\documentclass[10pt,a4paper]{amsart}
\usepackage[margin=19mm]{geometry}
\usepackage{amsmath,amssymb,mathtools}
\usepackage[scr=rsfs,cal=euler]{mathalfa}
\usepackage{xfrac}
\usepackage[unicode,hidelinks,bookmarks=false]{hyperref}
\newcommand{\ie}{i.e.\ }
\newcommand{\cf}{cf.\ }
\newcommand{\resp}{respectively\ }
\newcommand{\Subsection}{\subsection}
\providecommand{\nobreakdash}{\nobreak\hbox{-}}
\setlength{\emergencystretch}{2em}
\multlinegap0pt
\usepackage{amssymb}
\usepackage{xcolor}
\usepackage[shortlabels]{enumitem}
\usepackage[normalem]{ulem}
\usepackage{tikz}
\newtheorem{thm}{Theorem}
\newcommand{\idq}{{}^{\alpha}_{\omega}}
\title{Packing Independent Cliques in $K_4$-minor-free Graphs}
\author{Benjamin Xiao, Dong Ye}
\date{Source: arXiv:2512.18090v1 (CC BY 4.0)}
\begin{document}
\maketitle

\noindent\emph{Source and licence.} Packing Independent Cliques in $K_4$-minor-free Graphs. Version arXiv:2512.18090v1 (CC BY 4.0), distributed by arXiv under the Creative Commons Attribution 4.0 International licence, http://creativecommons.org/licenses/by/4.0/, as recorded in the arXiv licence metadata for this exact version and shown on the abstract page https://arxiv.org/abs/2512.18090v1. Full licence notice: \texttt{licenses/arxiv-2512.18090-CC-BY-4.0.txt}.\par\smallskip

\noindent\textit{Editorial scope.} Counted unit: the article's thm thm (source0.tex lines 767--771). The article's complete original proof of it is delivered with it (source0.tex lines 772--780, one \texttt{proof} environment of 9 lines). No mathematics is written, repaired or re-derived: the statement and the proof are the article's own lines, in the article's own notation, and no retained line is substituted or re-flowed.  No further source-local context is needed: every cross-reference of every retained line resolves inside the two delivered spans.  Nothing was omitted from any retained span, so the package carries no omission record.  No retained line contains a citation, so the wrapper carries no bibliography and no bibliography entry.  No retained line contains a figure environment, an image-inclusion command or drawing code, so no image is delivered and the package declares no asset dependency.  Editorial formatting only, in addition: an AMS \texttt{amsart} preamble that restates the article's own declarations and macros used by the retained lines, namely the environments \texttt{thm} on separate counters, together with the standard packages those lines need.  The retained environments print in this document as Theorem 1 in order of appearance, whereas the article prints the same items with the numbering of its own full document.  The complete native source of the article as distributed in the arXiv e-print for this exact version is bundled unchanged as \texttt{sources/191361/source0.tex}.
\begin{thm}
Let $G$ be an $n$-vertex subcubic graph. Then 
\[\idq(G)\ge \frac n 2,\]
and the bound is sharp. 
\end{thm}

\begin{proof}
    Let $G$ be a subcubic graph. Let $(X,Y)$ be a vertex partition of $G$ that maximizes the size of $E(X,Y)$, the set of edges joining a vertex of $X$ and a vertex of $Y$. If one part of $(X,Y)$, say $X$, has a vertex $v$ with degree at least two neighbors in its own part $X$. Then the partition $(X\backslash \{v\}, Y\cup \{v\})$ is a new vertex partition and 
    \[|E(X\backslash \{v\}, Y\cup \{v\})|>|E(X,Y)|,\] which contradicts that the maximality of $|E(X,Y)|$. Therefore, every vertex of $X$ or $Y$ has at most one neighbor in its own part, which means the vertex-induced subgraphs $G[X]$ and $G[Y]$ have maximum degree at most one. 
    Hence both $X$ and $Y$ are indeque sets of $G$. Note that $|X|+|Y|=|V(G)|=n$. It follows 
    \[\idq(G)\ge \max\{|X|, |Y|\}\ge \frac n 2.\]
    This completes the proof of the first part.
    
    To see the sharpness, there are 2-connected subcubic graphs with $4k$ vertices having indeque number $2k$: consider the subcubic graphs constructed from a cycle by blowing up every vertex to a square. It is quite easy to see that the graph has an indeque set of size $2k$. The graph does not have an indeque set of size at least $2k+1$. Otherwise, one of the 4-cycles contains at least three vertices from the indeque set, which is impossible because a 4-cycle has an indeque number two. 
\end{proof}

\end{document}
